\begin{exsol}{9}{Cabral 5.1 Front yards}{117}{Foundations}{B1}
\asked Judge whether the positive externality from a well-tended front yard is important, and
whether the market outcome is inefficient.

\method Externalities and market failure, Pigou corrections, social norms and the Coase theorem
(Cabral \S5.1, pp.~93--99); the Coase theorem and its failure conditions (\MR{} Ch.~2, pp.~35--39;
notes, Aula~2, ``The Coase theorem''). A two-neighbour model makes the comparison exact.

\solution
\subparagraph{(a) A two-sided externality: the model.} Two neighbours $i=1,2$ choose yard care
$x_i \ge 0$. Care gives the owner a benefit $a x_i$ (enjoyment, resale value), gives the
neighbour a spillover $s x_i$ with $s>0$, and costs $x_i^2/2$. Neighbour~1's payoff is
\[
u_1 = a x_1 + s x_2 - \tfrac12 x_1^2 ,
\]
and symmetrically for neighbour~2. The externality runs both ways: each is a producer and a
receiver of it.

\emph{Without bargaining} (each chooses alone), $\partial u_i/\partial x_i = a - x_i = 0$, so
$x_i^{N} = a$. The spillover $s$ does not enter the choice.

\emph{Efficient care} maximizes $u_1 + u_2$: $\partial(u_1+u_2)/\partial x_i = a + s - x_i = 0$,
so $x_i^{*} = a + s > a$. Each neighbour's gain from moving both to $x^*$ is
\[
u_1(x^*) - u_1(x^N) = \Big[(a+s)^2 - \tfrac12(a+s)^2\Big] - \Big[a^2 + sa - \tfrac12 a^2\Big]
= \tfrac12 s^2 ,
\]
so the joint loss of the uncoordinated outcome is $s^2$. With $a=1$, $s=\tfrac12$: care $1$
instead of $1.5$, and the pair loses $\tfrac14$. Without some mechanism, yards are
under-tended: the market outcome is inefficient, as with any positive externality (p.~94).

\subparagraph{(b) Is the externality important?} Per household the stakes are small: $s^2$ is
second order in the spillover, and a neighbour's yard is a modest part of the value of one's own
home. Two forces make it larger than it looks. A street has many neighbours, and the effects add
up across all who see the yard (the externality is close to a local public good, p.~95); and the
effect is capitalized into house prices, so a run-down block lowers every owner's wealth. Even
so, compared with pollution or congestion (pp.~94--95), it is a minor externality.

\subparagraph{(c) Can the neighbours bargain to the efficient level (Coase)?} With two
neighbours the Coase conditions nearly hold: few parties, a visible activity, low bargaining
costs, and no large wealth effects. The two-sided structure helps further: the deal is
reciprocal (``I keep my yard if you keep yours''), so with symmetric neighbours no money needs to
change hands; each gains $\tfrac12 s^2$. Two problems remain.
\begin{enumerate}[label=(\roman*)]
  \item \emph{Enforcement.} Given that the neighbour keeps $x_2 = a+s$, neighbour~1's best
  reply is still $x_1 = a$, which gains him $\tfrac12 s^2$ over keeping the deal. The agreement
  needs enforcement. Neighbours interact for years, so a social norm with the threat of
  withdrawn goodwill can sustain it (the repeated-game argument of p.~96).
  \item \emph{Many parties and measurement.} Along a whole street, each owner free rides on the
  others (the free-riding logic of p.~94), and ``well tended'' is hard to specify and verify;
  private information about how much each values tidiness makes a full agreement harder still
  (Aula~2, failure conditions of Coase).
\end{enumerate}
Where these costs bite, communities use other instruments: homeowners' associations with
covenants (a contract signed when the house is bought), municipal ordinances (direct regulation,
p.~96) and the social norm itself. A Pigou subsidy (p.~96) is possible in principle but
impractical because yard care is hard to measure.

\interp This is the Coase theorem at small scale. The externality creates an inefficiency only
to the extent that transaction costs block bargaining; in a two-household setting those costs
are low and norms do most of the work, so the residual inefficiency is small. On a large street,
free riding and measurement make explicit rules (covenants) the efficient organizational
response, which is \MR's corollary that efficient organization minimizes transaction costs
(p.~38).

\verify The sympy script \texttt{check\_B.py} confirms $x^N = a$, $x^* = a+s$, the gain
$\tfrac12 s^2$ per neighbour, the deviation gain $\tfrac12 s^2$ against the reciprocal deal, and
the loss $\tfrac14$ at $a=1$, $s=\tfrac12$. Limit case: $s \to 0$ gives $x^* \to x^N$ and no
loss; the inefficiency rises with the square of the spillover.

\answer{Yes, the outcome without bargaining is inefficient: each owner sets
$x^N = a < x^* = a+s$ because the neighbour's benefit is ignored, losing $\tfrac12 s^2$ per
household. The externality is modest per house but adds up across a street and is capitalized
into prices. Two neighbours can largely fix it by a reciprocal, norm-enforced deal (Coase); on a
whole street free riding and measurement push the solution towards covenants or ordinances.}
\end{exsol}
